\clearpage
\fancyhead[L]{\small PALLAS. CERES. VESTA.}
\begin{center}
\textbf{PALLAS. CERES. VESTA.}
\par\vspace{4pt}\hrule\vspace{4pt}
\fontsize{10}{12}\selectfont Göttingische gelehrte Anzeigen. Issue 136. Pp. 1348..1350. 26 August 1809.
\par\vspace{4pt}\hrule\vspace{4pt}
\end{center}

\par\indent Astronomers will likely welcome a brief report in our journal on the recovery of \emph{Pallas} at our observatory on the night of 9 to 10 August. The lack of usable observations from the previous year had kept Professor \textsc{Gauss} from improving the elements. Anyone looking for the planet must therefore still rely on the older elements, from which Professor \textsc{Bode} calculated a short ephemeris in the astronomical yearbook for 1811. The area where Pallas was expected was searched as early as 8 August, and the following evening one of the stars that had been marked down was found to have moved considerably south. On 10 August it was fully confirmed that this had been Pallas. Professor \textsc{Gauss} has so far made the following two observations [the same as those sent to von \textsc{Lindenau} in the letter of 14 Aug. 1809, p. 314 of this volume]. The first declination is somewhat less reliable than the second. Both right ascensions can be considered good. On this evidence, Professor \textsc{Bode}’s ephemeris represents the declination well but makes the right ascension 20 minutes too small. Pallas is already fully as bright as a star of the ninth magnitude.

\par\indent On 10 August \emph{Ceres} was also found again, and Professor \textsc{Gauss} compared it several times with nearby stars. These observations have yet to be fully reduced. For now it is enough to note that the planet was found within one minute of the position predicted by the XIII elements (see issue 87 of this year’s Gött. gel. Anz.) [3 June 1809, p. 313 of this volume] and by the ephemeris calculated from them (in the May issue of the Monatl. Correspondenz). Ceres also appears as bright as a star of the ninth magnitude.

\par\indent As early as 12 August Professor \textsc{Harding} also thought he recognized \emph{Vesta} in a faint little star of the ninth magnitude, very close to where it should be according to the IV elements (see issue 55 of these Anzeigen) [8 April 1809, p. 311 of this volume] and the ephemeris Dr. \textsc{Schumacher} calculated from them (Monatl. Corresp., at the place cited). Since then the weather has yet to allow any further confirmation.
